\begin{lemma}[Multi-sequence badness gives super-sequence badness]
Let $h$ be a bad locally constant multi-sequence, let $F^h$ be its
deciding front, and let $h^{\dcl}:F^h\to Q$ satisfy the value equation of
\noderef{DecidingValue}.  Then $h^{\dcl}$ is a bad super-sequence.
\end{lemma}
\begin{proof}
We prove the defining universal condition in
\noderef{BadSuperSequence}.  Fix finite sets $s,t\subseteq\mathbb N$,
membership proofs

\[
 h_s:s\in\mathsf{DecidingFrontSet}(h)
 \qquad\text{and}\qquad
 h_t:t\in\mathsf{DecidingFrontSet}(h),
\]

and assume $s\mathrel{\mathsf{FiniteShift}}t$.  By the definition of the
deciding front in \noderef{DecidingFrontSet}, these are precisely the
front-membership data needed to form the two subtype elements
$\langle s,h_s\rangle$ and $\langle t,h_t\rangle$ at which the value
field of $f$ is evaluated.

Unpack the finite-shift assumption using \noderef{FiniteShift}.  It gives
an increasing enumeration $x\in\mathsf{IncSeq}$ and two facts

\[
 p_s:\mathsf{ProperPrefixSet}(s,\operatorname{range}(x)),
 \qquad
 p_t:\mathsf{ProperPrefixSet}\bigl(t,
       \operatorname{range}(\mathsf{ShiftMap}(x))\bigr).
 \tag{1}
\]

Thus the witness is the same $x$ for both prefixes, and the second prefix
is taken in the shifted enumeration.  More explicitly, by the definition
of \noderef{ProperPrefixSet}, $p_s$ supplies some $n_s\in
\operatorname{range}(x)$ such that, for every $k$,

\[
 k\in s\quad\Longleftrightarrow\quad
 k\in\operatorname{range}(x)\ \text{ and }\ k<n_s,
\]

and $p_t$ supplies some $n_t\in
\operatorname{range}(\mathsf{ShiftMap}(x))$ such that, for every $k$,

\[
 k\in t\quad\Longleftrightarrow\quad
 k\in\operatorname{range}(\mathsf{ShiftMap}(x))\ \text{ and }\ k<n_t.
\]

The definition of \noderef{ShiftMap} makes $\mathsf{ShiftMap}(x)$ an
increasing enumeration, so both $x$ and $\mathsf{ShiftMap}(x)$ are valid
arguments for the supplied value-equation hypothesis $hf$.  Applying $hf$
to the first front element and the first prefix fact in (1) gives

\[
 e_s:\qquad
 f.value\langle s,h_s\rangle=h(x).
 \tag{2}
\]

Applying the same hypothesis separately to the second front element and
the second prefix fact in (1) gives

\[
 e_t:\qquad
 f.value\langle t,h_t\rangle=h(\mathsf{ShiftMap}(x)).
 \tag{3}
\]

These are applications of the theorem's particular hypothesis $hf$ to the
particular map $f.value$; the existential map supplied abstractly by
\noderef{DecidingValue} is not being substituted for $f$.

By \noderef{BadMultiSequence}, the badness hypothesis $hh$, instantiated
at this same enumeration $x$, gives

\[
 \neg r\bigl(h(x),h(\mathsf{ShiftMap}(x))\bigr).
 \tag{4}
\]

Suppose, toward a contradiction, that

\[
 r\bigl(f.value\langle s,h_s\rangle,
        f.value\langle t,h_t\rangle\bigr).
 \tag{5}
\]

Rewrite the first argument of (5) using (2) and the second using (3).
This yields $r(h(x),h(\mathsf{ShiftMap}(x)))$, contradicting (4).  Hence

\[
 \neg r\bigl(f.value\langle s,h_s\rangle,
             f.value\langle t,h_t\rangle\bigr).
\]

The finite sets, their front-membership proofs, and the finite-shift
assumption were arbitrary.  Therefore the displayed negation is exactly
the universal condition in \noderef{BadSuperSequence}, and $f$ is a bad
super-sequence.
\end{proof}
